\chapter{Electricity and Magnetism as One Relativistic Field}

We have studied how a prescribed electromagnetic field moves a charged
particle.  The unfinished half of the problem is the field itself: how
electric and magnetic fields evolve, and how charge and current act as their
sources.  Maxwell's equations answer those questions.  Before writing them,
however, it is worth recovering the physical puzzle that made their
relativistic structure unavoidable.

This is a deliberate retreat from formalism.  The action principle will
return, but not yet.  The point of beginning elsewhere is to see why the
electric and magnetic fields must be parts of one object before deriving
their dynamics from an action.

\section{The Puzzle in Einstein's First Paragraph}

During a visit to Belgium, a discussion of the early Solvay conferences led
back to Einstein's 1905 paper, \emph{On the Electrodynamics of Moving
Bodies}.  Its first paragraph states the motivation with remarkable
economy.  The same electromagnetic phenomenon receives two apparently
different explanations in two inertial frames.  One observer calls the
effect magnetic; another calls it electric.  The observed motion agrees, so
the two descriptions cannot represent two different pieces of physics.

The paper is now accessible at the level developed in this course.  Even if
one does not work through every calculation, the opening paragraph is worth
reading because it asks the right question before introducing the answer:
how can electric and magnetic descriptions change while the physical event
does not?

\subsection{A charge, a magnet, and two descriptions}

Take a magnetic field pointing in the \(y\)-direction and a charge moving
along the \(x\)-axis.  In the laboratory frame the electric field vanishes,
so the force is magnetic:
\begin{equation}
\mathbf F=q\,\mathbf v\times\mathbf B.
\end{equation}
It points in the remaining transverse direction.  The lecture's quick
blackboard axes are not meant to fix every right-hand-rule sign; the
important facts are that the force is perpendicular to both \(\mathbf v\)
and \(\mathbf B\), and that reversing the sign of the charge reverses the
force.

\begin{classroomqa}
\audiencequestion{Does the force point up or down in the drawing?}
\lecturerresponse{That depends on the sign of the charge.  Once the
orientation of the axes and the sign of \(q\) are fixed, the right-hand rule
settles the direction.  For the present argument only the existence of a
transverse force matters. \lecturetimestamp{00:07:13}}
\end{classroomqa}

Einstein described a conducting wire rather than an isolated particle.
Atoms provide an approximately fixed lattice while mobile charge carriers
move through it.  At the beginning of the twentieth century a current in a
wire was also much easier to observe than the trajectory of one free
charged particle.  The distinction is inessential here: a moving carrier in
a magnetic field feels the same Lorentz force.

\begin{figure}[H]
\centering
\includegraphics[width=\linewidth]{../../figures/lecture_08_figure_01.png}
\caption{The moving charge, transverse magnetic field, and magnet geometry
used to restate Einstein's electrodynamic puzzle.}
\par\smallskip
\textit{Lecture frame: 00:07:00.}
\end{figure}

Now keep the wire, or the test charge, fixed and move the magnet instead.
The changing magnetic configuration produces an electromotive force, which
in local language is an electric field.  That electric field pushes the
stationary charge in the same transverse direction.  Thus the two
descriptions are
\begin{align}
\text{moving charge, fixed magnet:}\quad&
q\,\mathbf v\times\mathbf B,
\\
\text{fixed charge, moving magnet:}\quad&
q\,\mathbf E'.
\end{align}
The relative motion is the same.  A theory of relativity must therefore
explain how a field that is purely magnetic in one frame acquires an
electric component in another.

Electricity and magnetism are not separate substances.  They are
frame-dependent parts of one relativistic field.  Einstein did not yet
possess the modern field-tensor notation, and he did not call the required
frame changes Lorentz transformations in this paragraph.  Nevertheless,
the transformation question is already present in full.

\section{The Electromagnetic Field Tensor}

The object that can carry electric and magnetic data together is the
antisymmetric tensor
\begin{equation}
F_{\mu\nu}
=
\partial_\mu A_\nu-\partial_\nu A_\mu.
\label{eq:field-tensor-definition}
\end{equation}
Greek indices run over \(0,1,2,3\), with \(x^0=t\); Latin indices run over
the three spatial values.  Antisymmetry,
\begin{equation}
F_{\mu\nu}=-F_{\nu\mu},
\end{equation}
sets the diagonal entries to zero and leaves six independent components.
Those six are precisely the three components of \(\mathbf E\) and the three
components of \(\mathbf B\).

With the mostly-plus metric
\begin{equation}
\eta_{\mu\nu}=\operatorname{diag}(-1,1,1,1),
\end{equation}
and the convention inherited from the charged-particle action in the
previous lecture, define
\begin{align}
E_m&=F_{0m}
=\partial_0A_m-\partial_mA_0,
\label{eq:e-from-a}
\\
B_i&=(\boldsymbol\nabla\times\mathbf A)_i.
\label{eq:b-from-a}
\end{align}
The lecture corrects a spatial index while writing the first line.  The
clean relation is the one above: the free spatial index must agree on both
sides.

\begin{figure}[H]
\centering
\includegraphics[width=\linewidth]{../../figures/lecture_08_figure_02.png}
\caption{The field-tensor definition and the mixed time--space component
identified with the electric field.}
\par\smallskip
\textit{Lecture frame: 00:13:15.}
\end{figure}

The magnetic relation is the spatial part of the same definition.  For
example,
\begin{equation}
B_x
=
\partial_yA_z-\partial_zA_y
=
F_{yz},
\end{equation}
with the other two components obtained by cyclic permutation.

\begin{figure}[H]
\centering
\includegraphics[width=\linewidth]{../../figures/lecture_08_figure_03.png}
\caption{The vector-potential formulas for the electric field and the
magnetic curl assembled on the board.}
\par\smallskip
\textit{Lecture frame: 00:14:00.}
\end{figure}

\subsection{The component array}

Let the first tensor index label rows and the second label columns, ordered
as \(0,x,y,z\).  Then
\begin{equation}
F_{\mu\nu}
=
\begin{pmatrix}
0&E_x&E_y&E_z\\
-E_x&0&B_z&-B_y\\
-E_y&-B_z&0&B_x\\
-E_z&B_y&-B_x&0
\end{pmatrix}.
\label{eq:f-lower-matrix}
\end{equation}
The mixed time--space edge contains \(\mathbf E\); the spatial
\(3\times3\) block contains \(\mathbf B\).  The lower triangle is not new
information, because it is the negative reflection of the upper triangle.

\begin{figure}[H]
\centering
\includegraphics[width=\linewidth]{../../figures/lecture_08_figure_04.png}
\caption{The completed lower-index electromagnetic tensor with its electric
edge and magnetic spatial block.}
\par\smallskip
\textit{Lecture frame: 00:17:15.}
\end{figure}

Raising both indices changes the sign of every electric entry and leaves
the magnetic entries unchanged:
\begin{equation}
F^{\mu\nu}
=
\begin{pmatrix}
0&-E_x&-E_y&-E_z\\
E_x&0&B_z&-B_y\\
E_y&-B_z&0&B_x\\
E_z&B_y&-B_x&0
\end{pmatrix}.
\label{eq:f-upper-matrix}
\end{equation}
An electric component carries one time index and therefore receives one
minus sign.  A magnetic component carries two spatial indices and receives
none.

\begin{classroomqa}
\audiencequestion{Am I reading the matrix correctly that the magnetic
components have no time index?}
\lecturerresponse{Yes.  They occupy the purely spatial block.  The electric
components are exactly the mixed time--space entries along the first row
and first column. \lecturetimestamp{00:18:55}}
\end{classroomqa}

\section{Boosting a Pure Magnetic Field}

Consider a boost along \(x\), with \(c=1\):
\begin{align}
t'&=\gamma(t-vx),
&
x'&=\gamma(x-vt),
&
y'&=y,
&
z'&=z,
\end{align}
where
\begin{equation}
\gamma=\frac{1}{\sqrt{1-v^2}}.
\end{equation}
In matrix form,
\begin{equation}
x'^\mu
=
\Lambda^\mu{}_\nu x^\nu,
\qquad
\Lambda^\mu{}_\nu
=
\begin{pmatrix}
\gamma&-\gamma v&0&0\\
-\gamma v&\gamma&0&0\\
0&0&1&0\\
0&0&0&1
\end{pmatrix}.
\label{eq:x-boost}
\end{equation}

\begin{classroomqa}
\audiencequestion{Is a minus sign missing in the boost matrix?}
\lecturerresponse{Yes.  The off-diagonal time--space entries carry the
minus signs required by \(t'=\gamma(t-vx)\) and
\(x'=\gamma(x-vt)\).  The correction is made before the field
transformation is used. \lecturetimestamp{00:23:50}}
\end{classroomqa}

\begin{figure}[H]
\centering
\includegraphics[width=\linewidth]{../../figures/lecture_08_figure_05.png}
\caption{The corrected Lorentz-boost matrix for motion along the
\(x\)-axis.}
\par\smallskip
\textit{Lecture frame: 00:24:00.}
\end{figure}

Every upper tensor index transforms like an upper vector index.  Therefore
\begin{equation}
F'^{\mu\nu}
=
\Lambda^\mu{}_\sigma
\Lambda^\nu{}_\tau
F^{\sigma\tau}.
\label{eq:f-boost-law}
\end{equation}
The repeated \(\sigma\) and \(\tau\) are dummy indices.  For a tensor with
more indices, one transformation matrix would appear for each index.

\begin{figure}[H]
\centering
\includegraphics[width=\linewidth]{../../figures/lecture_08_figure_06.png}
\caption{The rank-two Lorentz transformation law, with one boost matrix for
each field-tensor index.}
\par\smallskip
\textit{Lecture frame: 00:26:40.}
\end{figure}

\subsection{The component that becomes electric}

Assume the unprimed frame contains only a magnetic field \(B_y\).  In the
conventions of Equation~\eqref{eq:f-upper-matrix}, the relevant independent
component is
\begin{equation}
F^{xz}=-B_y.
\end{equation}
We seek the primed mixed component \(F'^{0z}\), because it carries the
primed \(z\)-component of the electric field.  Equation
\eqref{eq:f-boost-law} gives
\begin{align}
F'^{0z}
&=
\Lambda^0{}_\sigma
\Lambda^z{}_\tau F^{\sigma\tau}
\\
&=
\Lambda^0{}_x\Lambda^z{}_zF^{xz}
\\
&=
(-\gamma v)(1)(-B_y)
=
\gamma vB_y.
\end{align}
Since \(F'^{0z}=-E'_z\) in this convention,
\begin{equation}
E'_z=-\gamma vB_y.
\label{eq:boosted-electric-field}
\end{equation}
The sign depends on the declared axis and field conventions; the invariant
point of the calculation is
\begin{equation}
\lvert E'_z\rvert
=
\gamma\lvert vB_y\rvert
\ne0.
\end{equation}
The double sum collapses because the original field has only this one
independent magnetic component.

\begin{figure}[H]
\centering
\includegraphics[width=\linewidth]{../../figures/lecture_08_figure_07.png}
\caption{The selected tensor component after the boost, showing the factor
\(\gamma v\) multiplying the original magnetic component.}
\par\smallskip
\textit{Lecture frame: 00:32:30.}
\end{figure}

Choose the primed frame to move with the charge.  The charge is then at
rest, so there is no magnetic part \(q\,\mathbf v'\times\mathbf B'\) of its
force.  Yet the force does not disappear: it is \(q\mathbf E'\).  The
laboratory observer's magnetic force and the comoving observer's electric
force are two descriptions of the same event.

That is the answer to the puzzle in Einstein's first paragraph.  The
question occupied roughly the first forty minutes of this lecture; the
historical joke is that it occupied Einstein for much longer.  The
calculation is worth the time because it turns the slogan that electric and
magnetic fields mix into one explicit tensor component.

\section{Maxwell's Equations: Identities and Dynamics}

Now turn from the force on a charge to the field equations.  They will be
derived from an action in the next lecture.  Here the purpose is first to
state them, understand their two different logical roles, and put the
identity-based half into covariant form.

\begin{classroomqa}
\audiencequestion{Did we not already derive these equations?}
\lecturerresponse{We derived the equation of motion of a charged particle
in a prescribed electromagnetic field.  Maxwell's equations are equations
of motion for the field itself, including the response of the field to
charge and current. \lecturetimestamp{00:35:10}}
\end{classroomqa}

\begin{classroomqa}
\audiencequestion{How many Maxwell equations are there?}
\lecturerresponse{There are eight scalar component equations.  Vector
notation packages them as two scalar equations and two three-component
vector equations, conventionally called four Maxwell equations.
\lecturetimestamp{00:35:24}}
\end{classroomqa}

We continue to use units in which \(c=1\) and the usual vacuum constants
have been absorbed into the definitions of charge and current.  Four of
the eight component equations follow from the definition
\(F=dA\).  They are kinematic identities.  The other four are dynamical and
could change if the field action changed.

\subsection{Three elementary vector operations}

For a differentiable three-vector field \(\mathbf A\),
\begin{align}
(\boldsymbol\nabla\times\mathbf A)_x
&=
\partial_yA_z-\partial_zA_y,
\\
\boldsymbol\nabla\cdot\mathbf A
&=
\partial_xA_x+\partial_yA_y+\partial_zA_z.
\end{align}
The remaining curl components follow by cyclically permuting
\(x\to y\to z\to x\).  For a scalar field \(s\),
\begin{equation}
\boldsymbol\nabla s
=
(\partial_xs,\partial_ys,\partial_zs).
\end{equation}
Thus a curl is a vector, a divergence is a scalar, and a gradient is a
vector.

\begin{classroomqa}
\audiencequestion{Does anyone need the curl and divergence definitions
reviewed in full?}
\lecturerresponse{The component definitions are recorded, but the other
components need not all be expanded.  Cycling the coordinate labels
produces them. \lecturetimestamp{00:37:32}}
\end{classroomqa}

Commutativity of smooth partial derivatives implies two identities:
\begin{align}
\boldsymbol\nabla\cdot
(\boldsymbol\nabla\times\mathbf A)
&=0,
\label{eq:div-curl}
\\
\boldsymbol\nabla\times
(\boldsymbol\nabla s)
&=\mathbf0.
\label{eq:curl-grad}
\end{align}

\begin{classroomqa}
\audiencequestion{Should we write out the full component proof that the
divergence of a curl vanishes?}
\lecturerresponse{The class declines.  Expanding the components produces
pairs such as \(\partial_x\partial_yA_z\) and
\(-\partial_y\partial_xA_z\), which cancel because the derivative order may
be interchanged. \lecturetimestamp{00:40:13}}
\end{classroomqa}

\begin{figure}[H]
\centering
\includegraphics[width=\linewidth]{../../figures/lecture_08_figure_08.png}
\caption{The divergence-of-curl and curl-of-gradient identities used to
obtain half of Maxwell's equations.}
\par\smallskip
\textit{Lecture frame: 00:41:40.}
\end{figure}

\subsection{No magnetic source from the vector potential}

Taking the divergence of Equation~\eqref{eq:b-from-a} gives
\begin{equation}
\boldsymbol\nabla\cdot\mathbf B
=
\boldsymbol\nabla\cdot
(\boldsymbol\nabla\times\mathbf A)
=0.
\label{eq:div-b-zero}
\end{equation}
No equation of motion was used.  The result follows from the potential
description and smooth derivatives alone.

\begin{figure}[H]
\centering
\includegraphics[width=\linewidth]{../../figures/lecture_08_figure_09.png}
\caption{The divergence of the magnetic field reduced to the divergence of
a curl.}
\par\smallskip
\textit{Lecture frame: 00:43:05.}
\end{figure}

\begin{classroomqa}
\audiencequestion{What is the divergence of the electric field?}
\lecturerresponse{It is the electric charge density in the units being
used, \(\boldsymbol\nabla\cdot\mathbf E=\rho\).  A point charge is the
prototype of electric field lines spreading from a source.
\lecturetimestamp{00:43:20}}
\end{classroomqa}

At this elementary level, Equation~\eqref{eq:div-b-zero} says there is no
magnetic charge density.  A smooth, globally defined vector potential
cannot produce a magnetic monopole source.  Later one must qualify this
statement: monopoles can be described with potentials defined on
overlapping patches or with singular gauge choices.  The lecture postpones
that refinement.

\subsection{Faraday's equation from the electric potential}

Equation~\eqref{eq:e-from-a} is, in three-vector notation,
\begin{equation}
\mathbf E
=
\frac{\partial\mathbf A}{\partial t}
-\boldsymbol\nabla A_0.
\label{eq:e-vector-potential}
\end{equation}
Although \(A_0\) belongs to a four-vector, it is a scalar under ordinary
spatial rotations.  Take the curl:
\begin{align}
\boldsymbol\nabla\times\mathbf E
&=
\boldsymbol\nabla\times
\frac{\partial\mathbf A}{\partial t}
-
\boldsymbol\nabla\times
(\boldsymbol\nabla A_0)
\\
&=
\frac{\partial}{\partial t}
(\boldsymbol\nabla\times\mathbf A)
\\
&=
\frac{\partial\mathbf B}{\partial t}.
\end{align}
The spatial curl commutes with the time derivative, while the curl of the
gradient vanishes.

\begin{figure}[H]
\centering
\includegraphics[width=\linewidth]{../../figures/lecture_08_figure_10.png}
\caption{The magnetic curl, electric potential formula, and beginning of
the curl calculation in one completed board state.}
\par\smallskip
\textit{Lecture frame: 00:47:45.}
\end{figure}

The four homogeneous component equations are therefore
\begin{align}
\boldsymbol\nabla\cdot\mathbf B&=0,
\label{eq:homogeneous-div}
\\
\boldsymbol\nabla\times\mathbf E
-\frac{\partial\mathbf B}{\partial t}
&=\mathbf0.
\label{eq:homogeneous-curl}
\end{align}
The first is one scalar equation; the second is a three-component vector
equation.

\begin{figure}[H]
\centering
\includegraphics[width=\linewidth]{../../figures/lecture_08_figure_11.png}
\caption{The source-free Maxwell pair obtained directly from the
potential definitions.}
\par\smallskip
\textit{Lecture frame: 00:48:45.}
\end{figure}

\subsection{The dynamical half}

The remaining equations are
\begin{align}
\boldsymbol\nabla\cdot\mathbf E&=\rho,
\label{eq:gauss-electric}
\\
\boldsymbol\nabla\times\mathbf B
+\frac{\partial\mathbf E}{\partial t}
&=-\mathbf J.
\label{eq:ampere-maxwell}
\end{align}
They are not consequences of \(F=dA\).  They describe how sources create
and change the field, and the next lecture will derive them from the
electromagnetic action.  The minus sign on the current is fixed by the
conventions already chosen: \(F_{0i}=E_i\),
\(F_{ij}=\epsilon_{ijk}B_k\), and
\(\partial_\mu F^{\mu\nu}=J^\nu\).  The live board initially places
\(\mathbf J\) on the right without this minus sign; the resulting mismatch
with charge conservation is noticed and corrected in the following lecture.

\begin{classroomqa}
\audiencequestion{What kind of quantity is on the right-hand side of the
curl equation?}
\lecturerresponse{The left-hand side is a three-vector, so the right-hand
side is a three-vector called the electric current density
\(\mathbf J\).  The right-hand side of Gauss's equation is the
three-dimensional scalar charge density \(\rho\).
\lecturetimestamp{00:50:36}}
\end{classroomqa}

\begin{figure}[H]
\centering
\includegraphics[width=\linewidth]{../../figures/lecture_08_figure_12.png}
\caption{The four Maxwell equations grouped into the identity-based pair
and the sourced dynamical pair.}
\par\smallskip
\textit{Lecture frame: 00:50:40.}
\end{figure}

\section{Charge and Current as Spacetime Flow}

The symbols \(\rho\) and \(\mathbf J\) must now be given geometric meaning.
Charge density is charge per unit volume; current density is charge crossing
an oriented area per unit time.  The lecture develops both from worldlines
so that their later union as a four-vector will not look accidental.

\subsection{A classroom aside from the neighboring blackboard}

Before drawing those worldlines, equations visible in the neighboring room
occasion a brief teaching aside.  A forthcoming course would introduce
quantum mechanics to mathematics faculty who wanted the background needed
for quantum field theory and its connections to modern mathematics.  The
challenge runs in both directions: the physicist must make tacit physical
habits explicit, while the mathematician must learn what the formal
structures are being asked to describe.

\begin{classroomqa}
\audiencequestion{Who is expected to find that teaching assignment
frustrating?}
\lecturerresponse{Possibly both sides: physicists teaching an unfamiliar
audience and mathematicians confronting the physical habits behind the
formalism.  The assignment is nevertheless expected to be fun.  The class
then returns to charge and current. \lecturetimestamp{00:53:20}}
\end{classroomqa}

\subsection{Charge density}

Draw time vertically and let the remaining drawn directions stand for
three-dimensional space.  Charged particles trace worldlines.  At one fixed
time, place a small spatial cube of volume \(\Delta V\) across that bundle
of worldlines and add the charges of the particles inside it.  The charge
density is
\begin{equation}
\rho(t,\mathbf x)
=
\lim_{\Delta V\to0}
\frac{\Delta Q}{\Delta V}.
\label{eq:charge-density}
\end{equation}
The two-dimensional square on the board represents a three-dimensional
cube; the missing spatial direction is suppressed only because four
dimensions cannot be drawn on a flat board.

\begin{figure}[H]
\centering
\includegraphics[width=\linewidth]{../../figures/lecture_08_figure_13.png}
\caption{Charged-particle worldlines crossing a small spatial cell, with
\(\rho=\Delta Q/\Delta V\).}
\par\smallskip
\textit{Lecture frame: 00:57:20.}
\end{figure}

\subsection{Current density}

Turn the spacetime cell on its side.  To define \(J_x\), choose a small
window of area \(\Delta A_x\) in the \(yz\)-plane, perpendicular to the
\(x\)-axis.  Count the net charge that crosses from the negative-\(x\) side
to the positive-\(x\) side during \(\Delta t\):
\begin{equation}
J_x
=
\lim_{\Delta A_x,\Delta t\to0}
\frac{\Delta Q_x}{\Delta A_x\,\Delta t}.
\label{eq:jx-window}
\end{equation}
The orientation gives the sign.  Windows perpendicular to \(y\) and \(z\)
define \(J_y\) and \(J_z\) in the same way.

The board struggles honestly with the drawing: a window in the \(yz\)-plane
looks like a window in the ceiling rather than one in the wall.  There are
simply more relevant axes than can be displayed at once.  The invariant
definition is unambiguous even when the sketch is not.

\begin{classroomqa}
\audiencequestion{Is the window in the drawing on the ceiling?}
\lecturerresponse{In the room analogy, yes.  What matters is that its plane
is perpendicular to the \(x\)-direction.  A \(yz\) area measures the
\(x\)-directed current crossing it. \lecturetimestamp{00:59:38}}
\end{classroomqa}

\begin{figure}[H]
\centering
\includegraphics[width=\linewidth]{../../figures/lecture_08_figure_14.png}
\caption{The same worldline bundle viewed through a transverse window,
giving \(J_x=\Delta Q/(\Delta A\,\Delta t)\).}
\par\smallskip
\textit{Lecture frame: 01:01:40.}
\end{figure}

\begin{figure}[H]
\centering
\includegraphics[width=\linewidth]{../../figures/lecture_08_figure_15.png}
\caption{The three current components defined by charge crossing windows
normal to the three spatial axes.}
\par\smallskip
\textit{Lecture frame: 01:02:40.}
\end{figure}

The four quantities have parallel definitions.  A spatial volume at one
instant is a hypersurface perpendicular to the time axis, so \(\rho\) is
the flow of charge through a time-normal window.  The \(J_i\) are flows
through spatially normal windows.

\begin{classroomqa}
\audiencequestion{Are the four numbers one flow in the time direction and
three flows in the spatial directions?}
\lecturerresponse{That is exactly the useful picture.  A spatial volume is
a window perpendicular to the time axis; the three transverse areas are
windows perpendicular to \(x\), \(y\), and \(z\).  Their four fluxes belong
together. \lecturetimestamp{01:04:14}}
\end{classroomqa}

This suggests the contravariant four-current
\begin{equation}
J^\mu
=
(\rho,J_x,J_y,J_z).
\label{eq:four-current}
\end{equation}

\begin{classroomqa}
\audiencequestion{Where does the metric enter, and is this current
covariant or contravariant?}
\lecturerresponse{As defined here, \(\rho\) and the three \(J_i\) form the
components of \(J^\mu\), an upper-index or contravariant four-vector.  The
metric lowers it when \(J_\mu\) is required.  The covariant Maxwell
formulation will show why the upper-index choice is natural.
\lecturetimestamp{01:04:48}}
\end{classroomqa}

\section{Interlude at the Break: Signals Near a Black Hole}

At the scheduled break, questions turn to a previous discussion of black
holes.  Although this is a
digression from electrodynamics, it is a substantial classroom exchange
and illustrates the same discipline of separating descriptions made by
different observers.

\begin{classroomqa}
\audiencequestion{If a distant observer sees an infalling object slow near
the horizon, should the black hole look covered by everything that has ever
fallen onto it?}
\lecturerresponse{In Schwarzschild coordinate time, classical infall is
seen to approach the horizon asymptotically.  One can picture successive
matter as increasingly compressed layers near the horizon.  Operationally,
however, the outgoing signal is increasingly delayed, redshifted, and
faint.  Light emitted at an ordinary local frequency arrives first as
optical light, then infrared, microwave, and radio, with each photon
carrying less energy.  No realistic distant detector continues to see a
bright frozen object forever; it fades into blackness.
\lecturetimestamp{01:06:11}}
\end{classroomqa}

The infalling observer's clock and locally emitted flashlight remain
ordinary to that observer.  The distant description differs because the
frequency and arrival rate are redshifted.  With finite energy and finite
bandwidth, there is effectively a last resolvable part of the message.

\begin{classroomqa}
\audiencequestion{Is the object asymptotically approaching the event
horizon?}
\lecturerresponse{That is the distant Schwarzschild-coordinate
description.  Quantum mechanically the fading signal must also be
distinguished from the black hole's Hawking radiation.  At sufficiently
late times the weak redshifted signal is observationally lost in the
background of the hole's own radiation. \lecturetimestamp{01:09:24}}
\end{classroomqa}

\begin{classroomqa}
\audiencequestion{What if the redshifted wavelength grows beyond the scale
of the observable universe?}
\lecturerresponse{The signal is lost as a practical matter long before that
limit.  The lecture gives the order-of-magnitude Hawking picture of roughly
one typical quantum per light-crossing time,
\[
t_{\rm cross}\sim\frac{R_s}{c}.
\]
For a large hole this is extraordinarily weak radiation, and a sufficiently
redshifted infalling signal becomes undetectable much earlier than any
cosmological wavelength bound. \lecturetimestamp{01:09:50}}
\end{classroomqa}

\begin{classroomqa}
\audiencequestion{Does the rate of about one quantum per crossing time stay
fixed throughout the life of the black hole?}
\lecturerresponse{No.  An evaporating Schwarzschild black hole shrinks, so
its crossing time decreases and its Hawking temperature rises.  The
emission therefore accelerates.  The dimensional estimate in terms of the
instantaneous radius remains useful, but the radius is not constant.
Correspondingly, the typical Hawking wavelength shortens as evaporation
proceeds. \lecturetimestamp{01:11:42}}
\end{classroomqa}

After the break, with the jacket removed, we return to \(\rho\) and
\(\mathbf J\).

\section{Local Conservation of Charge}

It is tempting to state charge conservation only as
\begin{equation}
\frac{dQ_{\rm universe}}{dt}=0.
\end{equation}
That statement is too weak.  Charge might disappear from a laboratory and
reappear instantly at Alpha Centauri, or farther away in the Coma Cluster,
without changing the total.  Such a law would have little experimental
bite.

The stronger statement is local: charge inside a region can change only by
crossing the boundary of that region.  Place a small box around a point.
If its charge increases, there must be a net inward current through its six
faces; if it decreases, there must be a net outward current.

\subsection{Flux through opposite faces}

Take the two faces perpendicular to \(x\).  The inward contribution on one
face is evaluated at one value of \(x\); the contribution on the opposite
face is evaluated a small distance away.  Their difference, divided by the
box width, approaches
\begin{equation}
-\partial_xJ_x.
\end{equation}
The minus sign appears because positive \(J_x\) is defined as flow toward
increasing \(x\), while an outward flux decreases the charge in the box.
The \(y\)- and \(z\)-face pairs similarly contribute
\(-\partial_yJ_y\) and \(-\partial_zJ_z\).  Therefore
\begin{equation}
\frac{\partial\rho}{\partial t}
=
-\partial_xJ_x-\partial_yJ_y-\partial_zJ_z
=
-\boldsymbol\nabla\cdot\mathbf J.
\end{equation}
Equivalently,
\begin{equation}
\boxed{
\frac{\partial\rho}{\partial t}
+\boldsymbol\nabla\cdot\mathbf J=0
}.
\label{eq:continuity}
\end{equation}

\begin{figure}[H]
\centering
\includegraphics[width=\linewidth]{../../figures/lecture_08_figure_16.png}
\caption{The local balance law
\(\dot\rho=-\boldsymbol\nabla\cdot\mathbf J\), together with the small
box whose boundary currents produce it.}
\par\smallskip
\textit{Lecture frame: 01:20:45.}
\end{figure}

This is conventionally called the continuity equation.  The lecture
prefers the name \emph{local conservation}: continuity can suggest a
mathematical smoothness assumption, while the physical content is that
charge cannot disappear here and reappear there without a current in
between.

\subsection{Why Maxwell needed the displacement current}

Equation~\eqref{eq:continuity} is not an independent addition to Maxwell's
equations.  Take a time derivative of Gauss's law and a divergence of the
Ampere--Maxwell law:
\begin{align}
\frac{\partial\rho}{\partial t}
&=
\boldsymbol\nabla\cdot
\frac{\partial\mathbf E}{\partial t},
\\
-\boldsymbol\nabla\cdot\mathbf J
&=
\boldsymbol\nabla\cdot
(\boldsymbol\nabla\times\mathbf B)
+
\boldsymbol\nabla\cdot
\frac{\partial\mathbf E}{\partial t}
\\
&=
\boldsymbol\nabla\cdot
\frac{\partial\mathbf E}{\partial t}.
\end{align}
The two left-hand sides are therefore equal, giving
\(\partial_t\rho=-\boldsymbol\nabla\cdot\mathbf J\), exactly
Equation~\eqref{eq:continuity}.  What matters structurally is that the time
derivative of \(\mathbf E\), the displacement-current term, is essential for
the divergence relation to be compatible with charge conservation.

\begin{classroomqa}
\audiencequestion{We have not shown where the two sourced Maxwell equations
come from, have we?}
\lecturerresponse{No.  They are the dynamical target and will be derived
from the action next time.  The magnetostatic curl law was known before
Maxwell; Maxwell's crucial time-dependent addition was the displacement
current.  Its presence makes the field equation compatible with local
charge conservation. \lecturetimestamp{01:23:15}}
\end{classroomqa}

The detailed divergence calculation is left as an exercise in the lecture,
but its conclusion is not optional: Maxwell's equations require charge
conservation.

\section{The Four-Current and Its Divergence}

Write Equation~\eqref{eq:continuity} in components:
\begin{equation}
\frac{\partial\rho}{\partial t}
+
\frac{\partial J_x}{\partial x}
+
\frac{\partial J_y}{\partial y}
+
\frac{\partial J_z}{\partial z}
=0.
\label{eq:continuity-components}
\end{equation}

\begin{figure}[H]
\centering
\includegraphics[width=\linewidth]{../../figures/lecture_08_figure_17.png}
\caption{Local charge conservation expanded into one temporal and three
spatial derivatives.}
\par\smallskip
\textit{Lecture frame: 01:25:40.}
\end{figure}

With \(J^\mu\) defined by Equation~\eqref{eq:four-current}, the four terms
are one contraction:
\begin{equation}
\boxed{\partial_\mu J^\mu=0}.
\label{eq:four-current-conservation}
\end{equation}
The lower index on \(\partial_\mu\) contracts the upper index on
\(J^\mu\), so the result is a Lorentz scalar.  The geometric window
construction already suggested this transformation law: \(\rho\) and the
three current components are four orientations of the same charge flux.

\begin{figure}[H]
\centering
\includegraphics[width=\linewidth]{../../figures/lecture_08_figure_18.png}
\caption{The components \((\rho,J_x,J_y,J_z)\) assembled into \(J^\mu\)
and contracted with the four-gradient.}
\par\smallskip
\textit{Lecture frame: 01:27:10.}
\end{figure}

\begin{classroomqa}
\audiencequestion{When did the time derivative of \(\rho\) become a partial
derivative?}
\lecturerresponse{It was a partial derivative from the beginning.  We hold
the spatial box fixed and ask how its charge density changes as time
changes.  We are not following a moving trajectory, so
\(\dot\rho\) here means
\(\partial\rho/\partial t\) at fixed \(x,y,z\).
\lecturetimestamp{01:28:20}}
\end{classroomqa}

\section{The Homogeneous Equations in Covariant Form}

We now possess the required objects:
\begin{itemize}
\item \(\mathbf E\) and \(\mathbf B\) are the six independent components
of the antisymmetric tensor \(F_{\mu\nu}\);
\item \(\rho\) and \(\mathbf J\) are the four components of \(J^\mu\);
\item local conservation is the scalar equation
\(\partial_\mu J^\mu=0\).
\end{itemize}
The remaining task for this lecture is to write the homogeneous Maxwell
pair, Equations~\eqref{eq:homogeneous-div} and
\eqref{eq:homogeneous-curl}, as one tensor equation.

\subsection{The cyclic derivative}

The answer is
\begin{equation}
\boxed{
\partial_\sigma F_{\nu\tau}
+
\partial_\nu F_{\tau\sigma}
+
\partial_\tau F_{\sigma\nu}
=0
}.
\label{eq:bianchi}
\end{equation}
The three terms cyclically assign the derivative index and the two
field-tensor indices.  This is often abbreviated as
\begin{equation}
\partial_{[\sigma}F_{\nu\tau]}=0,
\end{equation}
where brackets denote antisymmetrization.

\begin{figure}[H]
\centering
\includegraphics[width=\linewidth]{../../figures/lecture_08_figure_19.png}
\caption{The cyclic three-index equation that packages the homogeneous
Maxwell equations.}
\par\smallskip
\textit{Lecture frame: 01:33:20.}
\end{figure}

\subsection{All three indices spatial}

Set
\begin{equation}
(\sigma,\nu,\tau)=(x,y,z).
\end{equation}
Then
\begin{align}
0
&=
\partial_xF_{yz}
+
\partial_yF_{zx}
+
\partial_zF_{xy}
\\
&=
\partial_xB_x+\partial_yB_y+\partial_zB_z
\\
&=
\boldsymbol\nabla\cdot\mathbf B.
\end{align}
Thus the all-spatial component of
Equation~\eqref{eq:bianchi} is precisely the absence of magnetic
divergence.

\begin{figure}[H]
\centering
\includegraphics[width=\linewidth]{../../figures/lecture_08_figure_20.png}
\caption{The all-spatial component of the cyclic identity reduced to
\(\partial_xB_x+\partial_yB_y+\partial_zB_z=0\).}
\par\smallskip
\textit{Lecture frame: 01:35:30.}
\end{figure}

\subsection{One time index}

Choose, for example,
\begin{equation}
(\sigma,\nu,\tau)=(y,x,0).
\end{equation}
Then Equation~\eqref{eq:bianchi} combines two spatial derivatives of
electric components with a time derivative of a magnetic component:
\begin{equation}
\partial_yF_{x0}
+
\partial_xF_{0y}
+
\partial_0F_{yx}
=0.
\end{equation}
Using the matrix in Equation~\eqref{eq:f-lower-matrix}, this is one
component of
\begin{equation}
\boldsymbol\nabla\times\mathbf E
-\frac{\partial\mathbf B}{\partial t}
=\mathbf0,
\end{equation}
with the relative signs fixed by antisymmetry.  Choosing the other pairs of
spatial indices supplies the other two curl components.

\begin{figure}[H]
\centering
\includegraphics[width=\linewidth]{../../figures/lecture_08_figure_21.png}
\caption{A one-time, two-space component check producing one component of
Faraday's equation.}
\par\smallskip
\textit{Lecture frame: 01:38:20.}
\end{figure}

There is also a direct proof.  Substitute
Equation~\eqref{eq:field-tensor-definition} into
Equation~\eqref{eq:bianchi}.  Six second-derivative terms appear, and they
cancel in pairs because partial derivatives commute.  This is the
four-dimensional version of the divergence-of-curl and curl-of-gradient
identities.

\begin{classroomqa}
\audiencequestion{Is there an \(h\) in the name Bianchi?}
\lecturerresponse{Yes.  The cyclic equation is the electromagnetic
Bianchi identity, a special case of a geometric identity that also appears
in general relativity.  Its deeper geometric interpretation is deferred.
\lecturetimestamp{01:40:28}}
\end{classroomqa}

\subsection{Why sixty-four written components contain only four equations}

Three indices, each with four possible values, appear to produce
\(4^3=64\) equations.  Antisymmetry removes the redundancy:
\begin{itemize}
\item if any two indices coincide, the cyclic expression vanishes
identically;
\item if all three indices are distinct, permutations reproduce the same
equation up to sign;
\item four dimensions contain only
\(\binom{4}{3}=4\) unordered triples of distinct indices.
\end{itemize}
There are \(24\) ordered triples with all indices distinct, grouped into
four sets of six permutations.  The remaining \(40\) ordered triples
vanish.  The four independent equations are one magnetic-divergence
equation and the three components of Faraday's equation.

\begin{classroomqa}
\audiencequestion{The field tensor has six independent entries, but the
three-index expression seems to contain sixty-four equations.  Why do only
four remain?}
\lecturerresponse{Repeated indices give zero, while permutations of three
different indices repeat the same relation up to sign.  The apparently
sixty-four components therefore collapse to four independent equations.
\lecturetimestamp{01:41:58}}
\end{classroomqa}

Equation~\eqref{eq:bianchi} is manifestly tensorial and therefore
Lorentz covariant.  In the smooth single-potential description it also says
there is neither magnetic charge density nor magnetic current.  The
dynamical sourced pair is left at exactly the boundary reached in the
lecture: its covariant equation and action derivation belong to the next
meeting.

\section{Duality and the Monopole Question}

The final questions return to the apparent asymmetry between electric and
magnetic sources.  In empty space, Maxwell's equations possess the
electric--magnetic duality rotation
\begin{equation}
\mathbf E\longrightarrow\mathbf B,
\qquad
\mathbf B\longrightarrow-\mathbf E.
\label{eq:duality-quarter-turn}
\end{equation}
This is a ninety-degree rotation in an abstract plane with axes
\(\mathbf E\) and \(\mathbf B\), not a reflection that merely interchanges
their labels.  With both electric and magnetic sources present, the
electric and magnetic four-currents must rotate at the same time.

\begin{figure}[H]
\centering
\includegraphics[width=\linewidth]{../../figures/lecture_08_figure_22.png}
\caption{The closing electric--magnetic duality picture: a quarter-turn in
the \(E\)-\(B\) plane rather than a simple interchange.}
\par\smallskip
\textit{Lecture frame: 01:45:00.}
\end{figure}

\begin{classroomqa}
\audiencequestion{If source-free equations are dual under electric and
magnetic interchange, why are the electric components time--space entries
while the magnetic components are space--space entries?}
\lecturerresponse{The distinction is frame-dependent within
\(F_{\mu\nu}\), and the source-free equations admit a duality rotation.
The correct operation is
\(\mathbf E\to\mathbf B\) and
\(\mathbf B\to-\mathbf E\), with source currents rotated similarly if
magnetic charge is included.  A direct interchange misses the required
sign. \lecturetimestamp{01:43:25}}
\end{classroomqa}

\begin{classroomqa}
\audiencequestion{Does the clean tensor picture make a magnetic monopole
meaningless?}
\lecturerresponse{No.  It means that a monopole cannot be represented by
one smooth vector potential over all of space.  The potential must instead
be patched, made singular in a gauge-dependent way, or replaced by an
equivalent dual description.  No elementary magnetic monopole had been
observed, and the lecture promises to return to how monopoles fit the
formalism rather than forcing the issue at the end of the night.
\lecturetimestamp{01:45:50}}
\end{classroomqa}

The lecture therefore ends at a precise frontier.  Relativity has explained
why electric and magnetic fields mix.  Potential definitions have supplied
the homogeneous Maxwell equations and their Bianchi form.  Charge flow has
become the four-current, and local conservation has become
\(\partial_\mu J^\mu=0\).  What remains is dynamical: the electromagnetic
action and the covariant sourced field equation.
